\begin{abstract}
Recent advancements, exemplified by BitNet~\cite{bitnet}, are ushering in an era where Large Language Models (LLMs) are effectively realized with only 1 bit per parameter. Here, we present a variant of the 1-bit LLM, denoted as \textbf{\text{BitNet b1.58}}, in which each model parameter assumes a ternary value from the set \{-1, 0, 1\}. Remarkably, this model achieves performance in terms of both perplexity and downstream task accuracy that is on par with full-precision Transformer LLMs (FP16 or BF16), provided equivalent model capacity and training data. At the same time, it offers substantial reductions in latency, memory footprint, throughput demands, and energy usage. More importantly, this 1.58-bit LLM proposes a novel scaling law and provides guidelines for constructing new high-efficiency LLM generations that do not compromise on capability. This direction also facilitates alternative computational strategies and lays the groundwork for developing custom hardware tailored for efficient 1-bit LLM processing.
\end{abstract}

\section{The Era of 1-bit LLMs}

In the last few years, the development of Large Language Models (LLMs) within artificial intelligence has accelerated, resulting in significantly larger and more powerful architectures. While these models have achieved outstanding success across diverse natural language processing tasks, their size brings obstacles regarding practical deployment as well as raising issues about their substantial energy usage and economic sustainability.

To mitigate these challenges, post-training quantization methods have been employed to generate low-bit representations of models, primarily for inference scenarios~\cite{smoothquant, gptq, quip, quip_sharp}. By reducing the precision of both weights and activations, these techniques substantially decrease the memory requirements and computational overhead associated with LLMs. Although there has been a progression from 16-bit to lower-precision implementations such as 4-bit variants~\cite{gptq, awq}, post-training quantization remains a less than ideal solution, even as it is commonplace in large-scale industrial deployments of LLMs.

Recent advances in 1-bit model designs, exemplified by BitNet~\cite{bitnet}, offer a compelling strategy for cutting LLM computational costs without sacrificing efficacy. Traditional LLMs typically operate with 16-bit floating point values (such as FP16 or BF16), with matrix multiplications constituting the primary workload. This places the main computational burden on floating-point arithmetic—specifically addition and multiplication. BitNet, in contrast, performs matrix multiplication using only integer addition, dramatically lowering energy expenditure for LLMs by several orders of magnitude. Since many hardware accelerators are power-limited, these energy savings also lead directly to increased processing speed.

Beyond computation, the movement of model weights from DRAM to accelerator-local memory (such as SRAM) is a substantial inference-time expense. Efforts to expand on-chip SRAM in order to boost throughput significantly raise hardware costs relative to DRAM. 1-bit LLMs present a considerably reduced memory demand, both in terms of capacity and bandwidth, when compared to models employing full-precision parameters. This leads to noticeable decreases in both time and cost required for weight loading during inference, promoting more efficient and rapid inference processes.

Here, we propose an advanced 1-bit LLM architecture termed \textbf{\text{BitNet b1.58}}, where each model parameter can take ternary values from the set \{-1, 0, 1\}. By introducing a zero value in addition to the original BitNet’s two-state configuration, \text{BitNet b1.58} operates effectively at 1.58 bits per parameter in binary encoding. This extension maintains all key advantages of 1-bit BitNet, such as its computationally novel framework that relies almost exclusively on addition operations for matrix multiplications, which can be extremely efficient. The energy consumption remains as low as that of the initial BitNet, and \text{BitNet b1.58} offers notable improvements in memory usage, throughput, and latency compared to FP16 LLM counterparts. Moreover, \text{BitNet b1.58} brings two notable improvements: its representational capacity is enhanced by allowing explicit feature filtering through zero weights, which helps to boost 1-bit LLM performance; additionally, our empirical results demonstrate that, starting at the 3B parameter scale and with identical settings (e.g., model size, number of training tokens), \text{BitNet b1.58} can reach parity with full-precision (FP16) models in both perplexity and task performance.

\section{BitNet b1.58}

\text{BitNet b1.58} derives from the BitNet framework, a Transformer variant in which the standard \emph{nn.Linear} operation is substituted with \emph{BitLinear}. The model is initialized from scratch, utilizing 1.58-bit quantized weights and activations quantized to 8 bits. Several changes from the initial BitNet model are introduced here, as detailed below.

\paragraph{Quantization Function.} Weight values are restricted to the set \{-1, 0, +1\} through the use of an \emph{absmean} quantization technique. This process involves scaling the entire weight matrix by the mean of its absolute values, and then mapping each entry to the closest integer within the target set:

\begin{equation}
    \widetilde{W} = \mathrm{RoundClip}(\frac{W}{\gamma + \epsilon}, -1, 1),
\end{equation}

\begin{equation}
    \mathrm{RoundClip}(x, a, b) = \max(a, \min(b, \mathrm{round}(x))),
\end{equation}

\begin{equation}
    \gamma = \frac{1}{nm}\sum_{ij} |W_{ij}|.
\end{equation}

The approach for quantizing activations maintains the methodology used in BitNet, except that we refrain from adjusting the activations into $[0, Q_b]$ prior to applying non-linearities. Rather, the activations are linearly mapped into the interval $[-Q_b, Q_b]$ for each token, removing the need for zero-point quantization. This approach not only simplifies the implementation and facilitates system-level optimizations, but also, according to our empirical findings, introduces minimal impact on overall model performance.

\paragraph{LLaMA-alike Components.}

LLaMA~\cite{llama, llama2} serves as a foundational architecture in many open-source large language models (LLMs). To align with prevailing open-source practices, \text{BitNet b1.58} is constructed with modules akin to those in LLaMA. In particular, the architecture incorporates RMSNorm~\cite{rmsnorm}, SwiGLU~\cite{swiglu}, and rotary embeddings~\cite{rope}, and eliminates the use of biases throughout. This compatibility allows \text{BitNet b1.58} to be integrated effortlessly with commonly used open-source frameworks such as Huggingface, vLLM~\cite{vllm}, and llama.cpp\footnote{\url{https://github.com/ggerganov/llama.cpp}}.

\section{Results}

We benchmarked \text{BitNet b1.58} against our implementation of the FP16 \text{LLaMA LLM} at multiple scales. To maintain parity across experiments, all models were pre-trained with 100 billion tokens using the RedPajama dataset~\cite{redpajama}. Their zero-shot abilities were then assessed on several language understanding benchmarks: ARC-Easy~\cite{arc}, ARC-Challenge~\cite{arc}, Hellaswag~\cite{hellaswag}, Winogrande~\cite{winoGrande}, PIQA~\cite{piqa}, OpenbookQA~\cite{openbookqa}, and BoolQ~\cite{boolq}. Additionally, we included validation perplexity evaluations on WikiText2~\cite{wikitext2} and C4~\cite{c4} to further analyze model performance.

To assess efficiency, we examined both runtime GPU memory utilization and inference latency for \text{LLaMA LLM} and \text{BitNet b1.58}, leveraging the FasterTransformer\footnote{\url{https://github.com/NVIDIA/FasterTransformer}} framework, which is highly optimized for low-latency LLM inference on GPUs. For \text{BitNet b1.58}, the 2-bit kernel from Ladder~\cite{ladder} was incorporated. Performance was quantified in terms of generation time per output token, as this represents the principal resource cost during inference.

A summary of the perplexity metrics and corresponding computational costs for \text{BitNet b1.58} and \text{LLaMA LLM} is provided in Table~\ref{tab:ppl}. The data illustrate that \text{BitNet b1.58} achieves parity in perplexity with full precision \text{LLaMA LLM} starting at the 3B parameter scale, while delivering a 2.71-fold improvement in speed and a 3.55-fold reduction in GPU memory demand. Notably, the 3.9B variant of \text{BitNet b1.58} attains a 2.4-fold speedup, requires 3.32 times less memory, and yet yields significantly better results than \text{LLaMA LLM} at 3B scale.

Table~\ref{tab:zero-shot} details the models’ zero-shot accuracy across downstream tasks. Evaluations adhered to the standardized process provided by \emph{lm-evaluation-harness}\footnote{\url{https://github.com/EleutherAI/lm-evaluation-harness}}. The findings indicate that the performance difference between \text{BitNet b1.58} and \text{LLaMA LLM} shrinks as model size increases. Critically, \text{BitNet b1.58} matches the full precision baseline from 3B parameters upwards. Consistent with the perplexity trends, the 3.9B \text{BitNet b1.58} surpasses the 3B \text{LLaMA LLM} model while incurring lower computational and memory overhead, thereby establishing \text{BitNet b1.58} as a Pareto improvement compared to leading large language models.

\paragraph{Memory and Latency}

We expanded our investigation to include larger models with 7B, 13B, and 70B parameters and analyzed the associated computational cost. As depicted in Figure~\ref{fig:latency-memory-energy}, both latency and memory usage exhibit favorable scaling properties, with observed speedup becoming more pronounced for larger model sizes. Notably, \text{BitNet b1.58} 70B demonstrates a 4.1-fold speed advantage relative to the \text{LLaMA LLM} benchmark. This improvement stems from the increasing computational expense of \emph{nn.Linear} operations as model size grows. A similar scaling effect is apparent in memory usage: since embeddings are stored in full precision, their relative share of memory consumption diminishes with larger models. For consistency, both latency and memory metrics were obtained using a 2-bit kernel, suggesting potential for additional gains through further optimization.

\paragraph{Energy}

In addition, we assessed the energy requirements associated with arithmetic operations for both \text{BitNet b1.58} and \text{LLaMA LLM}. Our analysis concentrated on matrix multiplications, given their dominance in the computational profile of LLMs. Figure~\ref{fig:energy} details the energy breakdown, revealing that the primary computation for \text{BitNet b1.58} involves INT8 additions, whereas \text{LLaMA LLM} predominantly relies on FP16 additions and multiplications. Based on the energy modeling from~\cite{energycost, pokebnn}, \text{BitNet b1.58} achieves a 71.4-fold reduction in energy expenditure for matrix multiplications on 7nm technology compared to its FP16 counterpart. Further, we present end-to-end energy consumption statistics for models processing sequences of 512 tokens. Our findings indicate that as model size increases, \text{BitNet b1.58} yields proportionally greater energy savings relative to the FP16 \text{LLaMA LLM}, which can be attributed to the expanding contribution of \emph{nn.Linear} operations, while overheads from other components lessen for larger architectures.

\paragraph{Throughput}

We conducted a throughput analysis comparing \text{BitNet b1.58} and \text{LLaMA LLM} with 70B parameters on a dual 80GB A100 GPU setup. Pipeline parallelism~\cite{gpipe} was employed to accommodate the \text{LLaMA LLM} 70B model on the hardware. Batch size was incrementally increased up to the maximum allowed by GPU memory, with sequence length fixed at 512. As shown in Table~\ref{tab:throughput}, \text{BitNet b1.58} 70B is capable of handling a batch size up to 11 times larger than that of \text{LLaMA LLM}, resulting in a throughput improvement by a factor of 8.9.

\textbf{\text{BitNet b1.58} establishes a novel scaling relationship between model capacity, performance, and inference efficiency.} Drawing upon Figures~\ref{fig:latency-memory-energy} and \ref{fig:energy}, we note the following equivalence points across different 1.58-bit and 16-bit model configurations:

\begin{itemize}
    \item The 13B BitNet b1.58 surpasses the 3B FP16 LLM regarding latency, memory consumption, and energy efficiency.
    \item The 30B BitNet b1.58 demonstrates better latency, memory, and energy characteristics than the 7B FP16 LLM.
    \item The 70B BitNet b1.58 outperforms the 13B FP16 LLM in terms of latency, memory footprint, and energy demand.
\end{itemize}

\paragraph{Training with 2T Tokens}

The total number of training tokens represents a significant variable in scaling LLM performance. To examine how \text{BitNet b1.58} handles extensive token counts, we trained a version with 2T tokens, adhering to the StableLM-3B~\cite{StableLM-3B-4E1T} data preparation strategy, which is currently the top open-source 3B model. Evaluation covered a benchmark suite consisting of Winogrande~\cite{winoGrande}, PIQA~\cite{piqa}, SciQ~\cite{sciq}, LAMBADA~\cite{lambada}, and ARC-easy~\cite{arc}. Table~\ref{tab:2t} presents the reported zero-shot accuracy; for tasks with both accuracy and normalized accuracy, their average was computed. StableLM 3b results at the 2T token scale are directly sourced from its technical report. Our analysis reveals that \text{BitNet b1.58} consistently exceeds StableLM’s performance across all evaluation tasks, underscoring the strong generalization abilities of 1.58-bit LLMs.

\section{Discussion and Future Work}

\textbf{1-bit Mixture-of-Experts (MoE) LLMs}

The Mixture-of-Experts (MoE) paradigm has emerged as an efficient method for implementing large language models (LLMs), offering significant computational savings in terms of FLOPs. However, practical use of MoE models is often constrained by substantial memory demands and the high costs associated with inter-chip communication. The advent of 1.58-bit LLMs presents promising solutions to these problems. With a notably smaller memory requirement, MoE architectures can be deployed with fewer hardware devices, reducing deployment complexity. Additionally, the reduction in memory use translates directly into lower communication costs when transferring activations across hardware, and, ideally, may enable fitting an entire model onto a single chip, thus eliminating network transfer overhead entirely.

\textbf{Native Support of Long Sequence in LLMs}

As language models increasingly require support for processing lengthy input sequences, managing the memory used by key-value (KV) caches becomes a prominent issue. \text{BitNet b1.58} makes important progress in this regard by reducing the bitwidth of activations from 16 bits to 8 bits, which allows for doubling the context window without extra resource consumption. Furthermore, future research can investigate methods to compress these activations even further, to 4 bits or less, especially for 1.58-bit LLMs, opening up avenues for even longer context support while maintaining efficiency.

\textbf{LLMs on Edge and Mobile}

Deploying 1.58-bit LLMs on edge and mobile platforms can lead to significant gains, particularly because such environments typically face tight constraints in terms of available memory and processing capability. Lowered memory and energy usage enable these compressed LLMs to be practical for a variety of tasks on devices where traditional models would be infeasible. This broadened applicability can substantially empower edge and mobile systems to leverage advanced language modeling. Additionally, the nature of 1.58-bit LLMs aligns well with the predominant use of CPU processors in these contexts, enhancing the execution efficiency of \text{BitNet b1.58} and further broadening the horizon of LLM deployments outside centralized cloud infrastructures.

\textbf{New Hardware for 1-bit LLMs}

Initiatives such as~Groq\footnote{\url{https://groq.com/}} have illustrated the feasibility and impact of constructing dedicated hardware accelerators (e.g., LPUs) specifically for LLM workloads. Building upon this foundation, we advocate for research and development efforts directed at novel hardware designs tailored for the unique computational requirements of 1-bit LLMs, as unlocked by the computational innovations of BitNet~\cite{bitnet}. 

\end{document}